\section{The Plan With One Fetch in It}
\label{sec:ch08-the-plan-with-one-fetch-in-it}

\index{query plan!asking for one|(}%
Ask the router to show its work by sending a header with the request:

\begin{minted}{text}
POST /graphql HTTP/1.1
Host: localhost:3002
Content-Type: application/json
X-WG-Include-Query-Plan: true

{"query":"{ sessions { nodes { id title speaker { name } } } }"}
\end{minted}

\index{ports!3002 is the only thing that changed}%
Note the \code{Host}. That is the only line separating this request from the
one chapter~\ref{ch:data-without-n-plus-1} sent to the service directly.

\index{query plan!the header does nothing at first}%
The response is correct, complete, and has no plan in it. No \code{extensions}
key, no error, no warning in the router's log. It is byte for byte what comes
back without the header.

\index{query plan!silence is the failure mode}%
The failure mode here is silence, and it is the only thing in this chapter that
behaves that way. A header this router will not honor and a
header it has never heard of produce the same response as each other, and that
response is the correct answer to the query you asked. There is nothing to
notice and nothing to search for.

\index{Advanced Request Tracing!gates the query plan}%
The reason is back in section~\ref{sec:ch08-a-router-with-nothing-to-serve}, in
a line I flagged and left. The query plan is an option of Advanced Request
Tracing, and ART is one of the five features the router listed as disabled when
it found no graph token. A plan tells the reader which of your internal services
holds what, and which query the router sent to each, so gating it behind
something is right. Gating it behind a token from a vendor's platform is a
choice, and locally it means one more line.

\subsection{The Whole File Again, With the Switch In}

\index{config.yaml@\code{config.yaml}!gains dev\_mode}%
\begin{minted}[highlightlines={13-17}]{yaml}
version: "1"

# Without this the router asks Cosmo Cloud for its execution config and
# refuses to start, because it has no token to ask with. The path is
# resolved against the directory the router is started in, not against
# this file, so start it from this folder:
#
#   router -config config.yaml
execution_config:
  file:
    path: ../graph/router.json

# A development switch, and what makes the query plan visible at all.
# The plan is an Advanced Request Tracing option, and outside dev mode
# the router wants a Cosmo Cloud token before it hands one out. With
# this off, X-WG-Include-Query-Plan is ignored rather than refused.
dev_mode: true
\end{minted}

\index{dev\_mode@\code{dev\_mode}!the router explains itself once it is on}%
Restart, and the router explains the mechanism to you in its own log without
being asked. The disabled-feature list drops to four names, Advanced Request
Tracing having left it, and two lines appear that were not there before:

\begin{minted}[fontsize=\footnotesize,breakanywhere]{text}
WARN  Development mode enabled. This should only be used for testing purposes
WARN  Advanced Request Tracing (ART) is enabled in development mode but requires a graph token to work in production. For more information see https://cosmo-docs.wundergraph.com/router/advanced-request-tracing-art
\end{minted}

\index{dev\_mode@\code{dev\_mode}!changes the log format too}%
Those two are trimmed the way the earlier blocks were, and there is more to
trim, because dev mode does something else nobody asks it for: the log stops
being JSON and becomes colored console output, each line gaining a timestamp
and the file and line in the router's own source that emitted it. Handy at a
terminal. Worth knowing before you point a log shipper at a router somebody
switched this on.

\subsection{Reading It}

\index{query plan!what one looks like}%
Send the same request again and the plan arrives in \code{extensions}:

\begin{minted}[fontsize=\footnotesize]{json}
{
  "queryPlan": {
    "version": "1",
    "kind": "Sequence",
    "children": [
      {
        "kind": "Single",
        "fetch": {
          "kind": "Single",
          "subgraphName": "sessions",
          "subgraphId": "0",
          "fetchId": 0,
          "query": "{\n    sessions {\n        nodes {\n            id\n            title\n            speaker {\n                name\n            }\n        }\n    }\n}"
        }
      }
    ],
    "normalizedQuery": "{sessions {nodes {id title speaker {name}}}}"
  }
}
\end{minted}

\index{query plan!Sequence and Single}%
A \code{Sequence} holding one child, which is a \code{Single} fetch. Those two
nouns are the whole vocabulary you need for now: a sequence is a list of steps
that happen in order, and a single is one request to one subgraph. This graph
has one subgraph, so the sequence has one step, and its \code{query} is the
query you sent.

\index{subgraphId@\code{subgraphId}!is the generated id from router.json}%
\code{subgraphName} is the name out of \code{graph.yaml}, and
\code{subgraphId} is \code{"0"}, which is the generated id
chapter~\ref{ch:composition} read out of the \code{subgraphs} list in the
execution config. Every piece of that plan traces back to a line in a file you
wrote or composed.

\index{query plan!has the same shape for every operation here}%
Try to make that shape change and it does not. I asked for two root
fields in one operation, for the mutation, and for the deprecated
\code{sessionById}. Sequence, one Single, one fetch, every time. There is one
service and nothing to decide, so the planner has no decision to record. What
chapter~\ref{ch:second-third-subgraph} adds is the second name in that list, and
from then on the plan is where you go to find out what the router decided to do
about it.

\subsection{Two Things the Plan Gives Away}

\index{router!rewrites literals into variables}%
The first was not what I was looking for. Send an inline argument:

\begin{minted}{graphql}
{ node(id: "U3BlYWtlcjox") { __typename } }
\end{minted}

and the \code{query} inside the plan is not that:

\begin{minted}{graphql}
query($a: ID!){
    node(id: $a){
        __typename
    }
}
\end{minted}

\index{router!normalization is what normalizedQuery records}%
The router lifted the literal into a variable before forwarding it. The mutation
gets the same treatment to its input object, and \code{sessionById(id: 1)}
arrives as \code{query(\$a: Int!)}. So what your subgraph receives is not what
your client sent, which matters the first time you read a subgraph's own log
looking for a query you recognize. \code{normalizedQuery} beside the fetch is
the router's record of the operation it actually planned. Why it does this I
cannot tell you; I can tell you that it does.

\index{introspection!is planned against a subgraph that does not exist}%
The second explains something from two sections ago. Ask for the plan of an
introspection query and the fetch names a subgraph that does not exist:

\begin{minted}[fontsize=\footnotesize]{json}
{
  "kind": "Single",
  "subgraphName": "introspection__schema&__type",
  "subgraphId": "introspection__schema&__type",
  "fetchId": 0
}
\end{minted}

\index{introspection!is answered by the router itself}%
No \code{query} key at all, because there is no request to send. The router
answers introspection out of the schema in its own execution config and never
asks your service. That is why the router could introspect happily with the
Sessions process stopped: it was never going to ask it.

\subsection{What the Router Costs This Graph}

\index{statement count!is unchanged by the router}%
Chapter~\ref{ch:data-without-n-plus-1}'s statement log is still running in the
service, so the question is answerable rather than arguable. The same query,
three ways:

\begin{center}
\begin{tabular}{lr}
\toprule
Request & Statements \\
\midrule
Direct to the subgraph on 5001 & 2 \\
Through the router on 3002 & 2 \\
Through the router, plan only & 0 \\
\bottomrule
\end{tabular}
\end{center}

\index{X-WG-Skip-Loader@\code{X-WG-Skip-Loader}!plans without executing}%
The third row is a second header, \code{X-WG-Skip-Loader: true}, which tells the
router to plan the operation and stop. It answers with \code{data} set to null
and the plan in \code{extensions}, and the service behind it does no work at
all. Useful, and a good habit: you can read a plan for an expensive query
without paying for it.

\index{router!adds nothing but a plan to a graph of one}%
The first two rows are the chapter's actual claim. A plan with one fetch in it
means the router does nothing to this graph except plan and forward, and the
two responses are identical. Figure~\ref{fig:ch08-one-fetch} is the two paths
beside each other.

\begin{figure}[htbp]
  \centering
  % One request, with and without the router: this figure fixes the book's
% router idiom the way ch01-one-field.tex fixed the timeline idiom itself.
% Every later chapter that draws the router copies what lane B draws here.
%
% Same idiom as figures/tikz/ch01-one-field.tex and
% figures/tikz/ch06-two-routes.tex: each lane is a timeline read left to
% right, ticks mark stages, and a brace under the stage that matters carries
% a one-line label. Axis, tick, event, span and spanlabel styles are copied
% from those two files unchanged.
%
% Ticks are not evenly spaced by rule; they are placed so that lane B's
% first tick (client posts to the router, port 3002) and last tick (2
% statements to SQLite) sit at the same x as lane A's first tick (client
% posts straight to Sessions, port 5001) and last tick, so the eye reads
% both rows as the same request. Lane A needs only 4 ticks and so gets wide,
% even gaps; lane B needs 6 ticks in the same span, so its gaps are tighter
% and every event label is pre-broken by hand (never left to automatic
% wrapping or hyphenation) to fit its column, the same discipline the two
% reference files use for every multi-line label.
%
% The asymmetry the figure argues: lane B's ticks 2 and 3, normalize and
% plan, do not exist in lane A at all. A brace under exactly those two ticks
% names them as the only work the router adds. A second, shorter brace
% under ticks 4 to 6 marks that from there on, lane B is doing the whole of
% lane A: post to Sessions, run resolvers, run the same 2 statements. The
% two braces sit side by side with a small gap between them and never
% overlap horizontally, so the figure carries both claims without crowding.
%
% No timing appears anywhere, by design: SPEC decision 19 prints no
% milliseconds in this book, so the horizontal axis is sequence only, and
% no lane is drawn narrower or wider in a way that would read as faster or
% slower.
%
% No quotation marks anywhere in this file. The book's strict Ascii mode
% (see check-chapter.psd1) makes the prose gate read a literal " as a quote
% character, while \" is LaTeX's diaeresis accent and fails to compile
% inside a node. Every identifier below is written bare, and the URL paths
% are broken onto their own line by hand rather than quoted or wrapped.
%
% Bare tikzpicture (house style): the figure environment, caption and label
% live at the call site in the section file.
\begin{tikzpicture}[
  x=1mm, y=1mm,
  every node/.append style={font=\sffamily\scriptsize},
  axis/.style={draw, line width=0.4pt,
    -{Straight Barb[length=1.4mm, width=1.6mm]}},
  tick/.style={draw, line width=0.4pt},
  event/.style={anchor=south, align=center, text width=24mm,
    inner sep=1pt},
  span/.style={draw, line width=0.4pt},
  spanlabel/.style={anchor=north, align=center, inner sep=2pt},
  lane/.style={anchor=west, font=\sffamily\scriptsize\itshape},
]

% ------------------------------------------------------------- lane A
\node[lane] at (0,13)
  {A. straight to Sessions, the way every earlier request went};

\draw[axis] (0,0) -- (150,0);
\foreach \x in {16,58,100,138}{
  \draw[tick] (\x,-1.5) -- (\x,1.5);
}

\node[event, text width=34mm] at (16,3)
  {client posts to\\localhost:5001/graphql};
\node[event, text width=28mm] at (58,3)
  {Sessions parses it,\\with no subgraph\\to choose};
\node[event, text width=22mm] at (100,3)
  {resolvers run};
\node[event, text width=24mm] at (138,3)
  {2 statements\\to SQLite};

% ------------------------------------------------------------- lane B
\begin{scope}[shift={(0,-40)}]
  \node[lane] at (0,20)
    {B. the same request, through the router};

  \draw[axis] (0,0) -- (150,0);
  \foreach \x in {16,42,69,94,116,138}{
    \draw[tick] (\x,-1.5) -- (\x,1.5);
  }

  \node[event, text width=24mm] at (16,3)
    {client posts to\\localhost:3002\\/graphql};
  \node[event, text width=26mm] at (42,3)
    {router normalizes\\the operation};
  \node[event, text width=25mm] at (69,3)
    {router plans:\\one fetch,\\subgraph sessions};
  \node[event, text width=22mm] at (94,3)
    {router posts to\\localhost:5001\\/graphql};
  \node[event, text width=20mm] at (116,3)
    {resolvers run};
  \node[event, text width=18mm] at (138,3)
    {2 statements\\to SQLite};

  \draw[span] (29,-4) -- (29,-7) -- (82,-7) -- (82,-4);
  \node[spanlabel] at (55,-8)
    {the only work the router adds,\\and it is the same two\\statements
     at the end};

  \draw[span] (84,-4) -- (84,-7) -- (147,-7) -- (147,-4);
  \node[spanlabel] at (115,-8)
    {this span is\\lane A, entire};
\end{scope}

\end{tikzpicture}

  \caption{The same query direct and through the router. The router adds two
  stages that did not exist before, normalizing the operation and planning it,
  and everything after them is the request chapter~\ref{ch:first-subgraph}
  already made. Both paths end in the same two statements, which is what a plan
  with one fetch in it costs.}
  \label{fig:ch08-one-fetch}
\end{figure}

\index{router!chapter 10 is where the count changes}%
Hold on to the number 2, because it is the last time it is free.
Chapter~\ref{ch:second-third-subgraph} puts a second service behind the same
router and the plan grows a step that fetches entities in bulk, and
chapter~\ref{ch:router-n-plus-1} is about what the resolver on the far side of
that step does with them. That is the difference between 2 and something much
worse.
\index{query plan!asking for one|)}%
